\begin{center}
    {\Huge\bfseries Квадратные трехчлены}
\end{center}

\problem{1}{%
На доске написаны девять приведённых квадратных трёхчленов
\[
  P_i(x)=x^2+a_i x+b_i,\qquad i=1,2,\ldots,9.
\]
Известно, что обе последовательности
$a_1,a_2,\ldots,a_9$ и $b_1,b_2,\ldots,b_9$ являются арифметическими
прогрессиями. Оказалось, что сумма всех девяти трёхчленов имеет хотя бы
один действительный корень. Какое наибольшее количество исходных
трёхчленов может не иметь действительных корней?}

% ИСТОЧНИК 2 (полностью): XIX Всероссийская математическая олимпиада школьников,
% заключительный этап, 1992--1993 учебный год, 9 класс, задача 427
% (для 10 класса задача 435). Автор: А. Перлин.
% Архивная карточка задачи с решениями:
% https://problems.ru/view_problem_details_new.php?id=109523
% Архив сборника заключительных этапов (условия и решения, задача 427):
% https://mathematikalpha.de/wp-content/uploads/2023/11/Allrussische_Mathematik-Olympiade.pdf
\problem{2}{%
Квадратный трёхчлен $f(x)$ разрешается заменить на один из трёхчленов
\[
  x^2 f\!\left(\frac1x+1\right)
  \qquad\text{или}\qquad
  (x-1)^2 f\!\left(\frac1{x-1}\right).
\]
Можно ли с помощью нескольких таких операций из трёхчлена
$x^2+4x+3$ получить трёхчлен $x^2+10x+9$?}

% ИСТОЧНИК 3 (полностью): XXXVII Всероссийская математическая олимпиада школьников,
% заключительный этап, 2010--2011 учебный год, 9 класс, второй день, задача 9.6.
% Авторы: И. И. Богданов, А. Гарбер. Официальный PDF (условия и решения):
% https://olympiads.mccme.ru/vmo/2011/final/day2.PDF
\problem{3}{%
У Пети и Коли в тетрадях записаны по два числа: первоначально у Пети
стоят $1$ и $2$, а у Коли --- $3$ и $4$. Раз в минуту каждый мальчик
составляет квадратный трёхчлен, корнями которого являются два числа из
его тетради. Если уравнение, полученное приравниванием этих двух
трёхчленов, имеет два различных корня, то один из мальчиков заменяет
свою пару чисел на эти корни; иначе ничего не происходит.

Какое второе число могло оказаться у Пети в тетради в тот момент, когда
первое стало равным $5$?}

% ИСТОЧНИК 4 (полностью): XXXII Всероссийская математическая олимпиада школьников,
% заключительный этап, 2005--2006 учебный год, 9 класс, второй день, задача 9.8
% (также дана 10 классу как задача 10.7). Автор: С. Л. Берлов.
% Официальный архивный PDF с условиями:
% https://olympiads.mccme.ru/vmo/32/sheets.pdf
% Карточка задачи с решением:
% https://problems.ru/view_problem_details_new.php?id=109857
\problem{4}{%
Дан квадратный трёхчлен
\[
  f(x)=x^2+ax+b.
\]
Уравнение $f(f(x))=0$ имеет четыре различных действительных корня,
сумма двух из которых равна $-1$. Докажите, что
\[
  b\le -\frac14.
\]}

% ИСТОЧНИК 5 (полностью): International Zhautykov Olympiad 2013,
% Day 2, Problem 1. Архив AoPS с полным комплектом задач:
% https://artofproblemsolving.com/downloads/printable_post_collections/3743.pdf
\problem{5}{%
Дан квадратный трёхчлен $p(x)$ с действительными коэффициентами.
Докажите, что существует натуральное число $n$, для которого уравнение
\[
  p(x)=\frac1n
\]
не имеет рациональных корней.}

% ИСТОЧНИК 6 (полностью): Korean Mathematical Olympiad 2017,
% national round, Day 1, Problem 4. Архив AoPS с полным комплектом задач:
% https://artofproblemsolving.com/downloads/printable_post_collections/559855
\problem{6}{%
Определим функцию $F\colon\mathbb R\to\mathbb R$ формулой
\[
F(x)=
\begin{cases}
\dfrac{1}{x-1}, & x>1,\\[4pt]
1, & x=1,\\[4pt]
\dfrac{x}{1-x}, & x<1.
\end{cases}
\]
Пусть $x_1$ --- положительное иррациональное число, являющееся корнем
квадратного многочлена с целыми коэффициентами, и
$x_{n+1}=F(x_n)$ для всех $n\ge1$. Докажите, что существуют различные
натуральные $k$ и $\ell$, для которых $x_k=x_\ell$.}

% ИСТОЧНИК 7 (полностью): Международная олимпиада «Туймаада», 2022,
% старшая лига, второй день, задача 5. Автор: Ф. Петров.
% Официальный PDF с условиями и решениями на английском языке:
% https://tuymaada.lensky-kray.ru/wp-content/uploads/2022/07/solutions-day2.pdf
% Русский комплект второго дня:
% https://tuymaada.lensky-kray.ru/wp-content/uploads/2022/08/Resheniya-2-den-corr.pdf
\problem{7}{%
Докажите, что приведённый квадратный трёхчлен
\[
  P(x)=x^2+ax+b,\qquad a,b\in\mathbb R,
\]
не может принимать в десяти последовательных целых точках значения,
каждое из которых является степенью двойки с целым неотрицательным
показателем.}

% ИСТОЧНИК 8 (полностью): 29th International Mathematical Olympiad,
% Canberra, Australia, 1988, Day 2, Problem 6.
% Официальный PDF IMO на английском языке:
% https://www.imo-official.org/assets/documents/problems/1988/1988_eng.pdf
% Официальная страница года:
% https://www.imo-official.org/problems/1988/
\problem{8}{%
Пусть $a$ и $b$ --- положительные целые числа, причём $ab+1$ делит
$a^2+b^2$. Докажите, что число
\[
  \frac{a^2+b^2}{ab+1}
\]
является полным квадратом.}